\begin{solapartat}
\punts{0,25} $E_\mathit{fotons} = hf = h\,\dfrac{c}{\lambda} = 3,98\times 10^{-19}\ \mathrm{J} = 2,48\ \mathrm{eV}$

Balanç d'energia:

\punts{0,5} $E_{C,\text{màx}} = hf + (-W_e)$

\punts{0,4} $E_{C,\text{màx}} = 2,48 - 1,20 = 1,28\ \mathrm{eV}$ \qquad \punts{0,1}

\destaca{Alternativament:}

\punts{0,15} $W_e = 1,20\ \mathrm{eV}\,\dfrac{1,602\times 10^{-19}\ \mathrm{J}}{1\ \mathrm{eV}} = 1,92\times 10^{-19}\ \mathrm{J}$

\punts{0,25} $E_{C,\text{màx}} = 3,98\times 10^{-19} - 1,92\times 10^{-19} = 2,06\times 10^{-19}\ \mathrm{J}$ \qquad \punts{0,1}
\end{solapartat}

\begin{solapartat}
\punts{1,25} L'energia cinètica dels electrons només depèn de la longitud d'ona de la llum incident, per tant, aquest paràmetre no canviarà.
L'únic efecte serà un augment del nombre d'electrons emesos, és a dir, augmentarà la intensitat de corrent que circula per la fotocèl·lula.
\end{solapartat}
